\begin{abstract}
    Il worm ILOVEYOU, diffusosi su scala globale nel maggio 2000, rappresenta uno degli esempi più significativi di attacco informatico basato sull'interazione tra propagazione tecnica e manipolazione psicologica degli utenti. 
    Una volta eseguito, il codice malevolo sfruttava Microsoft Outlook per inviare copie di sé ai contatti presenti nella rubrica della vittima, favorendo una rapida propagazione attraverso reti di relazioni personali e professionali. 
    Si esamina il caso considerando le tecniche di ingegneria sociale impiegate, le vulnerabilità dei sistemi informatici dell'epoca e le conseguenze sulla sicurezza dei dati.
    ILOVEYOU costituisce pertanto un caso di studio rilevante per comprendere l'evoluzione delle minacce informatiche e il ruolo del fattore umano nella cybersecurity.
\end{abstract}